% Generated draft: requires mathematical and source review.
\documentclass[11pt]{article}
\usepackage[margin=25mm]{geometry}
\usepackage{fontspec}
\setmainfont{Latin Modern Roman}
\usepackage{amsmath,amssymb,mathtools,unicode-math}
\setmathfont{Latin Modern Math}
\usepackage{enumitem,needspace,xurl,hyperref}
\setlength{\parindent}{0pt}
\setlength{\parskip}{5pt}
\setlength{\emergencystretch}{3em}
\newcommand{\N}{\mathbb{N}}
\newcommand{\Z}{\mathbb{Z}}
\newcommand{\Q}{\mathbb{Q}}
\newcommand{\R}{\mathbb{R}}
\newcommand{\C}{\mathbb{C}}
\newcommand{\F}{\mathbb{F}}
\newcommand{\A}{\mathbb{A}}
\newcommand{\PP}{\mathbb{P}}
\newcommand{\E}{\mathbb{E}}
\newcommand{\eps}{\varepsilon}
\newcommand{\dd}{\,\mathrm{d}}
\begin{document}
\section*{UnsolvedMath ID: 2100204}
\textbf{Source problem code:} catalogue entry
\subsection{Definitions and mathematical statement}
Consider a smooth integrable system with two commuting first integrals $F,H$ on a symplectic $4$-manifold, with momentum map $\Phi=(F,H)$. A rank-one critical point has $\operatorname{rank}d\Phi=1$. The problem restricts to the \emph{stable rank-one singularities described by V. V. Kalashnikov}; their defining normal forms, stability notion, and local or semiglobal neighborhood are not supplied. The source says that one action variable $I_1$ extends smoothly and the other $I_2$ is singular.

The first question is to describe the symplectic invariants of these singularities and decide whether the pair $(I_1,I_2)$ is a complete invariant for the intended symplectic classification. Equality of actions requires compatible cycle and normalization choices; the source does not spell out the allowed action-coordinate changes.

The source gives the following \emph{almost equivalent} formulation. On a common smooth neighborhood let $\omega_1,\omega_2$ be symplectic forms for which
\[
\{F,H\}_{\omega_1}=\{F,H\}_{\omega_2}=0,
\]
and suppose $(F,H)$ defines a Kalashnikov stable rank-one singularity for each form. Assume that their action variables agree. In the source's compatible primitive convention this means
\[
\oint_\gamma\alpha_1=\oint_\gamma\alpha_2
\quad\text{for every cycle $\gamma$ on every regular fiber }
\mathcal L_{f,h}=\{x:F(x)=f,\ H(x)=h\},
\qquad d\alpha_i=\omega_i\quad(i=1,2),
\]
where the $1$-forms are appropriately chosen on neighborhoods on which these comparisons are defined. Ask whether there is a smooth map $\psi$, on the intended neighborhood and across the singularity, satisfying
\[
F\circ\psi=F,\qquad H\circ\psi=H,\qquad
\psi^*\omega_2=\omega_1.
\]
The statement requests a smooth map; nondegeneracy would make it locally invertible, but a globally bijective neighborhood map is not explicitly required in the supplied wording. Exactness and compatible primitives, the comparison of cycles, and the size of the neighborhood must be fixed, rather than assumed to hold globally without justification.

The source notes that the simplest case is known. The supplied summary describes it as the one-degree-of-freedom case and reports no verified published complete answer for the general sufficiency question. It calls Kalashnikov's class non-degenerate, but supplies no definition establishing equivalence with any standard nondegeneracy notion; this extra label is not used to replace the source-specific class.
\subsection{Short English statement}
The question asks whether the smooth and singular action variables completely determine Kalashnikov's stable rank-one singularities up to symplectic equivalence. Its period formulation asks whether equal integrals on all regular-fiber cycles force a smooth change of coordinates preserving both integrals and relating the symplectic forms. The precise singularity class and neighborhood conventions are missing.
\subsection{Sources}
\begin{itemize}
\item Local UnsolvedMath registry, permanent ID 2100204, exact source code catalogue entry, ``Open Problems in Integrable Systems --- Topology of integrable systems, Lagrangian fibrations, and their invariants''; status open, classification OPEN-TRIAGE.
\item \emph{Open Problems, Questions, and Challenges in Finite-Dimensional Integrable Systems} (2018), Problem 2.4, \url{https://arxiv.org/abs/1804.03737}. The record identifies extraction from public arXiv source TeX, accessed July 29, 2026, and says the source presented the item as open and the simplest case as known.
\item The supplied research summary refers to V. V. Kalashnikov's work, Bolsinov--Fomenko's book, and ``semisimple'' invariant theory in the Russian integrable-systems literature. No precise titles, dates, theorem numbers, or definitions are supplied. The summary reports that it could not verify a published complete sufficiency theorem for the general rank-one case; this is an open/unverified assessment, not a new exhaustive literature conclusion.
\end{itemize}
\end{document}
